\usepackage{ctex}
\usepackage{geometry,graphicx,xcolor,color,subcaption}
\geometry{a4paper, top=2.5cm, bottom=2.5cm, left=2.8cm, right=2.8cm, headheight=14.5pt}
\usepackage{amssymb,amsmath,amsthm}                             % 数学字体
\usepackage{cancel}                                             % 数学划线（\cancel）
\usepackage{booktabs}   % 提供 \toprule, \midrule, \bottomrule    % 采用 Palatino 风格字体
\usepackage{multirow}   % 表格合并行（\multirow）
\usepackage{newclude,ulem}
\definecolor{winered}{rgb}{0.5,0,0}
\definecolor{structurecolor}{RGB}{24,44,82}     % 深藏青：章/节标题主色（与封面同系）
\definecolor{gold}{RGB}{168,132,74}             % 雅金：装饰线与点缀
\definecolor{main}{rgb}{0.5,0,0}
\definecolor{second}{RGB}{115,45,2}
\definecolor{third}{RGB}{0,80,80}
% --- 盒子配色（标签签颜色 + 同系浅底） ---
\definecolor{thmColor}{RGB}{24,44,82}     \definecolor{thmBg}{RGB}{240,243,249}
\definecolor{defColor}{RGB}{172,88,40}    \definecolor{defBg}{RGB}{253,248,238}
\definecolor{exColor}{RGB}{56,108,160}    \definecolor{exBg}{RGB}{240,246,252}
\definecolor{proColor}{RGB}{20,104,98}    \definecolor{proBg}{RGB}{236,247,246}
\definecolor{keyColor}{RGB}{148,111,55}   \definecolor{keyBg}{RGB}{250,245,235}
\definecolor{tipColor}{RGB}{46,120,70}    \definecolor{tipBg}{RGB}{238,247,240}
\definecolor{warnColor}{RGB}{188,64,54}   \definecolor{warnBg}{RGB}{252,240,238}
\usepackage[colorlinks,linkcolor = structurecolor]{hyperref}           % 定义引用的颜色
\usepackage{tcolorbox}
\tcbuselibrary{skins,breakable}
\usepackage{listings}
\usepackage{xcolor}
\lstdefinestyle{pythonstyle}{
    language=Python,
    basicstyle=\ttfamily\small,
    keywordstyle=\color{blue}\bfseries,
    commentstyle=\color{green!50!black}\itshape,
    stringstyle=\color{red},
    numberstyle=\tiny\color{gray},
    numbers=left,
    numbersep=8pt,
    frame=single,
    rulecolor=\color{black!30},
    breaklines=true,
    showstringspaces=false,
    tabsize=4,
    captionpos=b,
    xleftmargin=2em,
    framexleftmargin=1.5em
}
\lstset{style=pythonstyle}


% ------------------------------------------------------------%
% 定理/定义/例题环境：圆角彩色标题签 + 柔色底 + 柔和投影（tcolorbox）
\usepackage{etoolbox}
\newcounter{theorem}[chapter]
\renewcommand{\thetheorem}{\thechapter.\arabic{theorem}}

\tcbset{
    commonbox/.style={
        enhanced,
        coltitle=white,
        fonttitle=\heiti,
        attach boxed title to top left={xshift=5mm, yshift*=-\tcboxedtitleheight/2},
        boxed title style={
            frame hidden, opacityback=1,
            rounded corners, arc=4pt,
            top=1.2mm, bottom=1.2mm, left=2.5mm, right=2.5mm
        },
        boxrule=0pt, frame hidden,
        rounded corners, arc=9pt,
        drop fuzzy shadow=black!18!white,
        left=9mm, right=5mm, top=4.2mm, bottom=4.5mm,
        middle=1em,
        before skip=1.6em, after skip=1.1em,
        fontupper=\linespread{1.45}\selectfont
    }
}

% 定理类盒子：#1=额外选项 #2=名称 #3=标题注记
\newtcolorbox{thmBox}[3][]{commonbox,
    colback=thmBg, colbacktitle=thmColor,
    before title={\refstepcounter{theorem}},
    title={{#2}\ \thetheorem\ifstrempty{#3}{}{（#3）}}, #1}
\newenvironment{theorem}[1][]{\begin{thmBox}{定理}{#1}}{\end{thmBox}}
\newenvironment{lemma}[1][]{\begin{thmBox}{引理}{#1}}{\end{thmBox}}
\newenvironment{corollary}[1][]{\begin{thmBox}{推论}{#1}}{\end{thmBox}}

% 定义类盒子
\newtcolorbox{defBox}[3][]{commonbox,
    colback=defBg, colbacktitle=defColor,
    before title={\refstepcounter{theorem}},
    title={{#2}\ \thetheorem\ifstrempty{#3}{}{（#3）}}, #1}
\newenvironment{definition}[1][]{\begin{defBox}{定义}{#1}}{\end{defBox}}
\newenvironment{axiom}[1][]{\begin{defBox}{公理}{#1}}{\end{defBox}}

% 命题类盒子
\newtcolorbox{proBox}[3][]{commonbox,
    colback=proBg, colbacktitle=proColor,
    before title={\refstepcounter{theorem}},
    title={{#2}\ \thetheorem\ifstrempty{#3}{}{（#3）}}, #1}
\newenvironment{proposition}[1][]{\begin{proBox}{命题}{#1}}{\end{proBox}}
\newenvironment{property}[1][]{\begin{proBox}{性质}{#1}}{\end{proBox}}

% 例题类盒子（longexample 可跨页）
\newtcolorbox{exBox}[3][]{commonbox,
    colback=exBg, colbacktitle=exColor,
    before title={\refstepcounter{theorem}},
    title={{#2}\ \thetheorem\ifstrempty{#3}{}{（#3）}}, #1}
\newtcolorbox{longexBox}[3][]{breakable, commonbox,
    colback=exBg, colbacktitle=exColor,
    before title={\refstepcounter{theorem}},
    title={{#2}\ \thetheorem\ifstrempty{#3}{}{（#3）}}, #1}
\newenvironment{example}[1][]{\begin{exBox}{例题}{#1}}{\end{exBox}}
\newenvironment{longexample}[1][]{\begin{longexBox}{例题}{#1}}{\end{longexBox}}
\newenvironment{exercise}[1][]{\begin{exBox}{练习}{#1}}{\end{exBox}}

% 无编号强调盒：重点公式 / 技巧 / 注意
\newtcolorbox{keybox}[1][重点]{commonbox, colback=keyBg,  colbacktitle=keyColor,  title={#1}}
\newtcolorbox{tipbox}[1][技巧]{commonbox, colback=tipBg,  colbacktitle=tipColor,  title={#1}}
\newtcolorbox{warnbox}[1][注意]{commonbox, colback=warnBg, colbacktitle=warnColor, title={#1}}

% 证明 / 解 / 注 / 引言落款
\renewenvironment{proof}[1][推导]{\par\vspace{0.6em}\noindent{\heiti\color{structurecolor}#1：}\;\songti}{\hfill\textcolor{gold}{\rule{0.55em}{0.55em}}\par\vspace{0.3em}}
\newenvironment{solution}{\par\vspace{0.4em}\noindent\underline{{\heiti 解.}} \;\kaishu}{\par\vspace{0.4em}}
\newenvironment{remark}{\par\vspace{0.4em}\noindent\underline{{\heiti 注.}} \;\fangsong}{\par\vspace{0.4em}}
\newenvironment{note}{\par\vspace{0.4em}\noindent{\heiti\color{proColor}注：}\;}{\par\vspace{0.4em}}
\newcommand{\intro}[1]{\rightline{\parbox[t]{5cm}{\footnotesize \fangsong\quad\quad #1 }}}
% ------------------------------------------------------------%

\usepackage{enumerate}
\usepackage{enumitem}
\setlist[enumerate,1]{label=\color{structurecolor}\arabic*.}
\setlist[enumerate,2]{label=\color{structurecolor}(\arabic*).}
\setlist[enumerate,3]{label=\color{structurecolor}\Roman*.}
\setlist[enumerate,4]{label=\color{structurecolor}\Alph*.}
\usepackage{float}
\usepackage{tikz}
\usetikzlibrary{angles, quotes, arrows.meta, calc}
% 设置章形式
% ---------------------------------- %
% \setlength{\parindent}{0pt}
\usepackage{titlesec, titletoc}
\linespread{1.5}

\usepackage{fancyhdr}
\fancypagestyle{plain}{%
	\fancyhf{} % clear all header and footer fields
	\fancyfoot[C]{\small\color{structurecolor!80}\thepage}
	\renewcommand{\headrulewidth}{0pt}
	\renewcommand{\footrulewidth}{0pt}
}
% 正文页眉：章/节名 + 金色细线；页脚居中页码
\renewcommand{\chaptermark}[1]{\markboth{第\zhnumber{\thechapter}章\quad #1}{}}
\renewcommand{\sectionmark}[1]{\markright{§\thesection\quad #1}}
\fancypagestyle{mainstyle}{%
	\fancyhf{}
	\fancyhead[OC]{\small\kaishu\color{structurecolor!80}\rightmark}
	\fancyhead[EC]{\small\kaishu\color{structurecolor!80}\leftmark}
	\fancyfoot[C]{\small\color{structurecolor!80}\thepage}
	\renewcommand{\headrulewidth}{0.6pt}
	\renewcommand{\headrule}{{\color{gold}\hrule width\headwidth height 0.6pt}}
	\renewcommand{\footrulewidth}{0pt}
}

\titlecontents{chapter}[0em]{\addvspace{0.8em}}
	{\color{structurecolor}\large\heiti 第 \thecontentslabel 章\quad}{}{\hfill\color{structurecolor}\contentspage}
\titlecontents{section}[2.5em]{}
	{\thecontentslabel\quad}{}{\titlerule*[0.6pc]{ .}\;\color{structurecolor}\contentspage}

% 章标题：金色饰线夹「第 x 章」+ 特大黑体标题 + 金/灰双线收底
\titleformat{\chapter}[display]
	{\normalfont}
	{\centering{\color{gold}\rule[-0.5ex]{2.2em}{1.6pt}}\hspace{1.2em}{\color{structurecolor}\large\kaishu 第 \zhnumber{\arabic{chapter}} 章}\hspace{1.2em}{\color{gold}\rule[-0.5ex]{2.2em}{1.6pt}}}
	{1.6ex}
	{\color{structurecolor}\centering\Huge\heiti}
	[{\vspace{1.2ex}{\color{gold}\titlerule[1.2pt]}\vspace{-3.6pt}{\color{structurecolor!45}\titlerule[0.4pt]}}]

% 节标题：金色 § 编号 + 黑体 + 金/灰双细线
\titleformat{\section}[hang]
	{\normalfont\Large\heiti\color{structurecolor}}
	{{\color{gold}\Large\S}\,\thesection\hspace{1em}}{0pt}{}
	[{\vspace{0.4ex}{\color{gold!85!black}\titlerule[0.9pt]}\vspace{-2.8pt}{\color{structurecolor!35}\titlerule[0.4pt]}}]

\titlespacing*{\section}{1pc}{*7}{*2.3}[1pc]
% 小节：金色短杠 + 深藏青黑体
\titleformat{\subsection}[hang]
	{\normalfont\large\heiti\color{structurecolor}}
	{{\color{gold}\rule[-0.3ex]{0.85em}{2.2pt}}\hspace{0.6em}\thesubsection\hspace{0.6em}}{0pt}{}
% 子小节：金色方块 + 黑体
\titleformat{\subsubsection}[hang]
	{\normalfont\normalsize\heiti\color{structurecolor!90}}
	{{\color{gold}\small$\blacksquare$}\hspace{0.5em}\thesubsubsection\hspace{0.6em}}{0pt}{}

% 与主题统一的通用盒子（标题由用户输入，供自定义题型）
% 导言区定义格式
\newtcolorbox{myTheorem}[1]{commonbox,
    colback=thmBg, colbacktitle=thmColor, title={#1}}
\newtcolorbox{myhypothesis}[1]{commonbox,
    colback=tipBg, colbacktitle=tipColor, title={#1}}
\newtcolorbox{mydefinition}[1]{commonbox,
    colback=defBg, colbacktitle=defColor, title={#1}}

% 公式强调框（无标题，居中）
\tcbuselibrary{skins, breakable, theorems}
\newtcolorbox{myformula}[1][]{
    enhanced, frame hidden, boxrule=0pt,
    colback=keyBg,
    rounded corners, arc=9pt,
    drop fuzzy shadow=black!18!white,
    before skip=1.2em, after skip=1.2em,
    halign=center, valign=center,
    top=3mm, bottom=3mm, left=4mm, right=4mm,
    #1
}
\definecolor{DarkBlue}{RGB}{0,0,139}

% ============================================================
%  工程数学特有设置（保留原主文件中的自定义命令）
% ============================================================
\usepackage{esint}          % \oiint 等面积分符号

% 自定义数学命令
\DeclareMathOperator{\re}{Re}
\DeclareMathOperator{\im}{Im}
\DeclareMathOperator{\Arg}{Arg}
\newcommand{\argz}{\operatorname{Arg}}   % 兼容原文档中的 \argz
\newcommand{\Ln}{\operatorname{Ln}}
\newcommand{\dif}{\mathop{}\!\mathrm{d}}

% 全局书签数学命令处理
\pdfstringdefDisableCommands{%
	\def\re{Re}\def\im{Im}\def\Arg{Arg}\def\argz{Arg}%
	\def\Ln{Ln}\def\dif{d}%
}